\section{HIV and the basic model of viral replication}

In the 1980s there was a major global HIV epidemic. This was terrible, and the
disease has killed on the order of forty million people since, but it had the
side effect of inducing a flurry of research into HIV pathogenesis and viral
dynamics more generally. By the early 1990s progress was being made with the
development of effective mathematical models describing within host viral
replication.

Principal among these was what became known as the Basic Model of Viral
Replication (BMVR)
\cite{mclean1993balance,perelson1996hiv,nowak1996population}, which emerged
from the fray of model diversity following a period of substantial
experimentation. Its basic assumption was quite simple: there are three agent
types.

\begin{itemize}
\item \textbf{Target cells: $T$}

Host cells produced with a constant influx of rate $s$, removed as they become
infected by individual virions with a rate $\gamma$ per cell--virion pair.

\item \textbf{Free virions: $V$}

Produced at a constant rate, $p$, by each infected cell, and also cleared with
rate $c$.

\item \textbf{Infected cells: $I$}

A single free virion coming into contact with a target cell is sufficient to
convert it into an infected cell, which removes one cell from $T$. Infected
cells then produce new free virions at a constant rate, not accounting for the
likely reality that a higher intracellular load might make a cell more
productive, though some models do include a lag period between the time a cell
is infected and when it begins to produce virions. They are also removed by
cell death with rate $d_I$.
\end{itemize}

\noindent
The equations are as follows,

\begin{align*}
\ddt{T} & = s - \gamma T V\\
\ddt{I} & = \gamma T V - d_I I\\
\ddt{V}& = p I - c V.
\end{align*}

\noindent
In some cases the model is simplified by taking the number of target cells as a
constant, $T$. This represents the situation in which there is a large
availability of cells to infect, such that their depletion is not a significant
driver of the overall dynamics. With this simplifying assumption the equations
become

\begin{align*}
\ddt{I} & = \gamma T V - d_I I\\
\ddt{V}& = p I - c V.
\end{align*}

\noindent
This is the form carried through
\mextref{p:ch:one-cell-to-population}{Chapter~6}. To be clear, it rests on the
assumption that virions are produced by cells one by one at a constant release
rate, and that release and cell death are separate, independent events.

\subsubsection*{New research into HIV viral release}

In 2019 an experiment was published \cite{hataye2019principles} with the
intention of quantifying the release
dynamics of HIV virions from infected cells. Small wells were each seeded with
a single infected cell, and viral counts were then measured from the liquid of
every well at intervals over a period of days.

If the assumption of the BMVR were correct, we would expect to see something
approximating a Poisson process in each well, with roughly the same count
recorded in all of the wells at each interval. What was actually observed was
very different: huge variation between wells, and production rates that were
not at all constant. Instead, a high degree of intermittency was recorded ---
cells remaining completely inactive for successive intervals, then occasionally
releasing a substantial number of virions, seemingly all at once. This
challenges the standard assumption of constant-rate budding directly, at the
level of the single cell where the assumption is made.

\subsubsection*{So what?}

The observation begs the question: why? There will be some reason HIV has
evolved to use this quasi-bursting procedure. Perhaps it is simply more
convenient for release to take place in floods. Passage through the cell wall
requires an opening, and it surely makes sense to induce this only
periodically, avoiding both unwanted leakage of cellular components and ingress
of extracellular debris.

But perhaps there is also some general advantage to be gained by virions being
released in large cohorts like this. It could be that locally depletable immune
factors are flooded and overcome by the sudden outflux of enemy agents. A paper
by N.~Komarova \cite{komarova2007viral} presented this idea of a so-called
``flooding'' advantage, and it has an obvious analogy: a gang surrounded in a
house may reason that if they all run out at once, some will get away in the
confusion. Whether the analogy survives contact with a mathematical model is
the subject of \mextref{path:ch:pathogenesis}{Chapter~8}, with attention paid
to the spectrum between budding release and complete cell bursting.

Back on the topic of the BMVR and its questionable assumption, the two
constants $p$ and $d_I$ are phenomenological. They are fitted to viral load
data and then reported as though they characterised the pathogen, when in fact
nothing in the model says what the cell is doing between infection and release.
For a budding virus that costs nothing, because release and death really are
independent events. For a lytic pathogen it is wrong in mechanism: such a
pathogen releases nothing while it replicates and then releases everything at
once, in the event that ends the cell's life. Release and death are not
independent; they are the same event.
\Mextref{p:ch:one-cell-to-population}{Chapter~6} replaces the two constants
with two functions of cell age derived from the birth--death--catastrophe
process, and recovers the BMVR as a projection of the result.
